% =====================================================================
%  c00.tex —— 标题页 / 译者说明 / 目次（Lead）
% =====================================================================
\thispagestyle{empty}
\begin{center}
\vspace*{2.2cm}
{\Huge\bfseries 数物系のための圏論}\\[0.6cm]
{\Large\bfseries ——导出范畴、三角范畴、$A_\infty$ 范畴}\\[1.6cm]
{\large 梶浦 宏成\ \ 著}\\[2.4cm]
{\large 中\ 文\ 译\ 注\ 本}\\[0.8cm]
{\small（据サイエンス社 SGC ライブラリ 75，2010 年初版译出）}
\vfill
{\small 中译与批注：本译本\\[0.4em]
2026 年}
\end{center}
\clearpage

\chapter*{前言}
\addcontentsline{toc}{chapter}{前言}
\markboth{前言}{前言}

本书是一本关于范畴论之同调代数侧面的入门书，特别解说导出范畴、三角范畴、$A_\infty$ 范畴以及它们所具有的性质。阅读时所需的知识几乎没有（本来是这样打算的）。唯独线性代数，希望读者有一定程度的熟悉。

「我是某某」——只要报上名字，即使不看对方的脸也能知道是谁。即使是陌生人，只要名字不同，也能知道是别人。若能再报出血型与生日就更好了\footnote{当然不可说谎。无矛盾地定义了的数学操作不会说谎。}！所谓同调就是这样的东西。它指的是把若干数列之类对应于数学对象的办法，给出分类乍看难以区分的数学对象的手段。而为了确定同调而建立的代数理论，一般称为同调代数。在几何中，同调例如对球面（球的表面）与环面（甜甜圈的表面）分别指派 $(1,0,1)$、$(1,2,1)$ 这样一组数。也就是说球面与环面是彼此不同的东西。然而，同调代数并不只能对付几何对象。那么在怎样的状况下可以使用同调代数？要思考这一点，就必须把「当时需要的是哪些性质」在头脑中整理清楚。适合做这种头脑整理的东西，就是范畴。

提起范畴（圏），人们会联想到「大气圈」「首都圈」之类的词，而在数学中，范畴也大致意味着所考虑的数学对象的范围。然而它不仅仅指定范围，还要看属于该范围之物之间的相关关系等结构。也就是说，范畴是数学对象之族的「容器」，把东西放进这个容器，事物就被整理好了。关于容器本身的一般理论也许显得枯燥无味，但一旦放入东西，就会觉察到它的妙处。同调代数也是如此：只要对象能放进某种性质足够好的容器（即后文所说的阿贝尔范畴），就能以共通的方式展开讨论。实际上，范畴这一概念，似乎正是以这种同调代数的形式化为发端的\footnote{作者不是适合谈论历史的人。参见 [38]、[39] 等。}。在范畴论中，同调代数之王般君临的，就是导出范畴，或者描述导出范畴结构的三角范畴。导出范畴是把作为同调之素材的东西范畴化而得的，实际上其范畴结构甚至包含了超出同调的信息。（原始文献为 [P.~Deligne, \textit{Cohomologie \'etale}, S\'eminaire de g\'eom\'etrie alg\'ebrique du Bois-Marie SGA 4 1/2, Lecture Notes in Mathematics 569, Springer.]，为法语。）

$A_\infty$ 范畴是较这些东西更新的概念，但近年来与弦理论相关的数理物理相联系而被热烈讨论，且与三角范畴很合得来。这是因为三角范畴所具有的超出同调的信息与 $A_\infty$ 结构有关；遗憾的是本书未能说明到其细节，但在正文的「三角 $A_\infty$ 范畴」或「三角范畴的 $A_\infty$ 增广」这一概念中，可以看到两者理论的融合。本书的构成，就是以通往这个三角 $A_\infty$ 范畴的道路为中心轴而安排的。

总之由于这样的缘由，本书选取导出范畴、三角范畴、$A_\infty$ 范畴这三者为中心人物。在物理学中为弦理论作范畴论式表述时，起最中心作用的也将是这三者。此外还有若干更有趣、且至今仍在发展的东西，但依作者偏颇的主观，本书没有涉及它们。例如模型范畴（Quillen [48]，cf.\ 土屋 [55]）、高阶范畴（它与 $A_\infty$ 范畴的高阶积这一意义没有直接关系，cf.\ [39, 第 XII 章]）等；后者应用于拓扑场论的数学表述的课题也正在被讨论（但作者对这一方向并不持肯定态度）。此外，与弦理论关系密切的三角范畴的稳定性条件（Bridgeland，cf.\ [1]）等也未能涉及。

本书所讨论内容的数学与物理背景，已在开头第 1 章作了一通说明，但那里出现了大量未加定义的专门术语。暂且把它当作各关键词之间相关关系的表格来读即可。从第 2 章起，将从基础的东西开始逐一给出定义。导来范畴、三角范畴这类形式而沉重的工具是本书的主题，本书采取的方針是：以尽量少的知识尽可能单刀直入地定义它们，进而应用于以下三角 $2\times2$ 矩阵及其表示为中心的具体例子，由此品味它们。也就是说，为了例子而包含了环论与表示论的讨论，但本书并不足以习得这些领域本身的标准基础。各章末给出了相关参考文献的清单，请一并参照。此外，例如在《数理科学 2008 年 3 月号：现代物理的数学关键词》中，有专家对相关概念的解说，特别是小林正典氏的「范畴与函子」、浅芝秀人氏的「Quiver 的表示」等与本书直接相关。又，《数理科学 2008 年 9 月号：拓扑弦的世界》中包含了大量与物理背景（1.4 节）相关的话题。另外，本书所涉及的基本事项大多在岩波书店的《数学辞典》中也有解说，不过目前还没有关于 $A_\infty$ 代数的记述。

（这可以说是关于数学书的一般性建议）阅读本书时有应当注意的一点。首先，暂且不要在意定理、引理等的证明细节，先把握话题的大框架，这是高效理解的诀窍。而当定义过于抽象、难以理解时，看例子有助于理解。当定理的陈述过于抽象、难以理解时，看它的证明有时也有助于理解，但基本上可以放到后面。实际上，从开头起就一路追着证明读本书是危险的。本书是一次尝试：不假定任何知识，铺一条通到今后研究有望进展之处的道路；但其中多少有些勉强，难度并不均匀。特别是关于 $A_\infty$ 结构的若干基本定理，本打算不省略证明地写出，却在那之前就达到（超过了）页数上限，这是遗憾之处。不过实际上，或许这样正好。总体而言，写着「证明概略」的地方，恐怕难以完全理解。写着「细节略」的地方，只是省略了占版面的、用矩阵进行的说明等，并不那么难。

在本稿的执笔过程中，谨向耐心引导总想继续写下去的作者直到脱稿的、サイエンス社的平势耕介氏表示感谢。

\begin{flushright}
2010 年 3 月\\[0.3em]
梶浦宏成
\end{flushright}

\clearpage

\chapter*{译者说明}
\addcontentsline{toc}{chapter}{译者说明}
\markboth{译者说明}{译者说明}

本译本是梶浦宏成《数物系のための圏論 —— 導来圏，三角圏，A∞ 圏を中心に》（サイエンス社，SGC ライブラリ 75，2010 年
7 月 25 日初版）的中文译注本，仅供学术研读之用。

本译本的处理原则如下：

\begin{enumerate}[leftmargin=2em,itemsep=3pt,topsep=4pt]
\item \textbf{公式与编号。}书中编号公式（如式 (3.8)）与定理类标题（如定義 3.29、定理 5.74）的编号一律照抄原书，
      便于与原文逐条对照；原书中大量未编号的公式保持不编号。
\item \textbf{正誤表。}作者曾发布《SGC75 正誤表》（2014 年 5 月 20 日版，另有 2014 年 5 月 24 日版关于八面体公理
      (Tr5) 的专项订正，原书 PDF 末尾附有）。凡正誤表涉及的页面，正文按订正后的内容翻译，并在相应位置以译者批注
      注明底本原印内容与订正依据。
\item \textbf{批注。}以「Tutor Note」形式（LaTeX 注释，纸面上不显示）给出：正誤表条目的对照、原文省略的论证、
      可验证的数值或符号锚点、术语的英文原词与译名依据。全部批注另汇编为《译注汇总》。
\item \textbf{术语。}中文有稳定通行译名者用中文（范畴、函子、态射、模、复形、正合列等），无通行译名者保留英文
      （operad、DG、$A_\infty$、Fukaya 等）；人名一律保留拉丁原名。术语对照见附录。
\item \textbf{图表。}原书的交换图表、箭图（箙）等一律用 TikZ 重画，节点标签译为中文，结构与原图逐项对应。
\end{enumerate}

底本页码与原书印刷页码一致（引用处标注为「原文 p.N」）。

\clearpage
\tableofcontents
\clearpage
